% Section paragraphs written during the image check; citations identify typescript pages.
\subsection*{14. Homotopical algebra and the limits of present topology (pp. 1--2)}
Grothendieck opens the continuation of the letter to Quillen by distinguishing the study of spaces of maps and their successive homotopies from the study of spaces of homeomorphisms, immersions, embeddings and fibrations. His terms ``wildness'', ``isotopy'', ``stratified structures'' and ``tubular neighbourhoods'' delimit problems for which the familiar category of topological spaces seems poorly adapted; the contrast with piecewise linear structures is that their maps are too sparse, notably because real-valued piecewise linear functions are not closed under multiplication. This develops the complaint recalled explicitly from the letter, p.~1, and supplies the contrast on which section~15 builds its more favorable account of localization and homotopy types. The interlinear revisions narrow the opening description to ``maps between spaces'' and replace ``numerical'' by $\mathbb R$-valued (p.~2). The witness differs in two mathematically relevant places: the typescript says ``pro-topoi'' on p.~1, where arXiv has ``pretopoi'', and arXiv retains ``numerical'' without applying its handwritten replacement on p.~2. These readings are established by the scan; the edition does not identify the two notions or silently repair either authorial formulation.

\subsection*{15. Models of homotopy types and the place of stacks (pp. 2--4)}
The argument passes from the failure of the Hauptvermutung as a literal assertion about triangulations to its much weaker homotopical analogue: appropriate localizations of topological spaces and semi-simplicial sets yield the same homotopy category. Grothendieck compares CW-spaces, Kan complexes and Serre's cubical complexes, then proposes stacks, previously called ``$\infty$-groupoïdes'', as a third description particularly close to the homotopy groups themselves. His analogy with chain complexes and the Dold--Puppe equivalence makes the intended role of noncommutative structure explicit; section~16 immediately generalizes the motivation from spaces to topoi. The distinction between describing and computing invariants matters: on p.~3 he twice strikes ``description'' and writes ``computation'', corrections which the arXiv witness does not apply. On p.~4 a marginal note dated 3.5 begins ``This is an error'' and retracts the claim that stack cohomology has no immediate description from the primitive structure, referring instead to the nerve and to forgetting structure. Broad pencil strokes across this note make its final cancellation status uncertain, and one verb remains doubtful; both uncertainties are preserved. The arXiv footnote is substantially longer than this marginal wording and must not be substituted for it as though it were a literal transcription.

\subsection*{16. Stacks over topoi and noncommutative cohomological algebra (pp. 4--6)}
The motivation for stacks is enlarged from section~15's models of homotopy types to coefficients for cohomology over an arbitrary topos. Starting with Giraud's $0$-, $1$- and $2$-dimensional theory, Grothendieck treats $n$-stacks as coefficients whose ``integration'' is evaluation at the final object, with the homotopy groups of the resulting stack supplying cohomology after a dimension shift. His letters to Breen and Mme Raynaud's work on Lefschetz comparisons provide the geometric motivation: higher noncommutative cohomology should unify homological and homotopical comparison statements. The passage to infinite-order stacks is presented as a qualitative enlargement of perspective, and the following section asks how much of the homotopy-theoretic side ordinary categories can already model. The scans matter for the formulas: p.~4 has $H^i$, not the witness's $K^i$; p.~5 has $R\Gamma F_*$ and complexes $F_*$, where the supplied arXiv comparison reports a missing $F$ and a $j$ in place of the star. The handwritten qualification ``[non commutative]'' on p.~5 must also remain visible. These are witness differences, not authorial mathematical corrections introduced by this edition.

\subsection*{16$'$. Categories and the ``impressive bunch'' of models (pp. 7--9)}
Grothendieck associates a homotopy type to a small category both through its presheaf topos and through its nerve, comparing the first construction with Verdier and Artin--Mazur pro-homotopy types. He recalls and rejects his former guess that passage to the enveloping groupoid preserves this type: Quillen had pointed out that ordinary categories realize arbitrary homotopy types. This corrects the initial groupoid intuition while continuing sections~15--16's search for models, and motivates the later study of the ``impressive bunch'' and of coherators. A second proposed shortcut is itself explicitly retracted: the margin of p.~7, dated 5.3, says that the asserted weak equivalence $K\longrightarrow NN'(K)$ is false and directs the reader to section~24, pp.~21--22. The note is shorter than its arXiv counterpart; its call number is (18), and the section has been renumbered $16'$ rather than left as $16$bis. Page~8 proposes cofibrations injective on objects and arrows with a further call (19), then contrasts general $n$-stacks with $n$-Gr-stacks and their truncated homotopy types. Its wording ``left derived functors'' for the inverse-limit cohomology construction should be treated as an authorial formulation requiring care, not silently changed to the usual right-derived terminology. The final paragraph on p.~9 introduces the dependence on a coherator and the universal-algebra question, despite professing little inclination to pursue it.

\subsection*{17. The Artin--Mazur cohomological criterion (p. 9)}
This short section supplements the unification program by recalling that weak equivalences can be detected through cohomology with twisted coefficients, including noncommutative $H^0$ and $H^1$ as well as the commutative higher groups. Grothendieck explains how this permits étale cohomological information to constrain homotopy types, while wondering whether its success discouraged the development of higher noncommutative cohomology. The reference back to section~15 is the author's own handwritten cross-reference and is retained even though the immediate unification discussion occupies section~16 in the present numbering. The section bridges that motivation and the concrete corrections of section~18. Its handwritten insertion reads ``(namely, unification)'', not the arXiv ``messy unification''; the final sentence deletes ``that'' and says ``their observation'', distinctions that are obscured by the witness. The section number itself is corrected from 18 to 17 on the scan.


\subsection*{18. Corrections to the coherator and higher coherence (pp. 10--11)}
The return to the letter (L) is explicitly corrective. Grothendieck withdraws the existence of arbitrary binary amalgamated sums in $\underline B_o$: the geometric intuition of gluing contractible cells along embedded subcells had concealed the presence of nonmonomorphic maps, including the codiagonal. He also revises the construction of $\underline C_{n+1}$ from $\underline C_n$ by insisting on the total set of new arrows, and rejects the assertion that intersections of maximal subcells must be subcells. The marginal drawing on p.~10 makes the latter failure concrete: three regions $f,f',f''$ meet so that $f\cap f''=\{s,t\}$. The scan therefore supplies a two-point intersection, not the quite different formula in the initial transcription. A separate correction explains why associativity beyond primary composition cannot be represented by a single homotopy arrow: the sources and targets are themselves only homotopic, so squares, cubes and higher coherence data intervene. Page~11 names the dual categories left and right coherators, then connects the proposed non-groupoidal stacks to Benabou's bicategories and their cellular, topological interpretation. This prepares section~19's realization and singular-stack adjunction. The source wording $L\to K$ and $M\to K$ beside $K\amalg_L M$ on p.~10 is retained as an apparent authorial slip; it does not give the usual span for that pushout. Calls (11) and (14) on p.~11 are present without the explanatory footnotes printed in arXiv.

\subsection*{19. A provisional itinerary and the sheaf temptation (pp. 12--13)}
The first two items of the itinerary are the definition of $\underline B_o$, its dual and the coherators, followed by a realization/singular-stack adjunction with topological spaces. Unlike arbitrary presheaves on a left coherator, Gr-stacks must carry the designated amalgamated sums to fiber products. The realization functor nevertheless extends to all presheaves by preserving colimits; its values on basic cells are prescribed, while extensions to the coherator require choices. The anticipated independence of those choices is therefore tied to an as-yet undeveloped notion of equivalence between Gr-stacks. On p.~13 Grothendieck asks whether the exactness conditions could instead be sheaf conditions for a topology, which would make Gr-stacks a topos, and proposes weak equivalences detected by realization, monomorphisms as cofibrations and a lifting-property definition of fibrations. These are programmatic proposals rather than established results. The ``quick glance'' suggesting a topology is retracted immediately in section~20; even the concluding sentence about the inclusion, revised to composite, functor is already struck on p.~13. The source has ``one reward'', not the witness's ``one ward'', and handwritten ``as cofibrations'', not ``for cofibrations''. Its call (13) contains no footnote text on this page.

\subsection*{20. Model categories, sites and a retracted canonical topology (pp. 14--15)}
The section opens by correcting section~19's ``quick glance'': the proposed topology on the coherator would make the component functors fiber functors and force componentwise colimits, which the author recognizes as false. The category of categories supplies a warning and a new candidate construction. Grothendieck considers families of functors whose nerves cover every finite composable string, proposes identifying them with universally effectively epimorphic families, and associates to a category $A$ the topos of sheaves on $(\mathrm{Cat})_{/A}$. The handwritten margin immediately marks the asserted equivalence as false and points to section~24; the digits are readable, while pencil hatching leaves the note's cancellation status uncertain. Page~15 broadens the proposal to a natural site structure on model categories, with the homotopy type of $X$ identified with that of the induced site or topos over $X$. Its further presumption that the topology should be canonical is itself contradicted by a dated margin note absent from the arXiv comparison. These local corrections are essential to reading the section as an evolving investigation. The four arrows in the diagram of effective, universal and ordinary epimorphisms are handwritten, as are the topos tildes in $(M_{/X})^{\sim}=M^{\sim}_{/X}$; the object labels are typed. The inquiry leads into section~21's wider inventory of models rather than establishing a general theorem about canonical sites.

\subsection*{21. Multicomplexes, symmetry and restricted coherence (pp. 15--16)}
The inventory now includes symmetric variants of semi-simplicial and cubical complexes and their multiple-complex analogues. A diagonal complex already describes the intended homotopy type, but Grothendieck explicitly pursues the variety of algebraic presentations, not a unique most convenient model. He compares the passage from complexes to multicomplexes with that from ordinary categories to higher categories or stacks, recalling Quillen's suggestion of defining higher categories through multiple semi-simplicial complexes. The return to the coherator is concrete: the coherence diagrams examined in Benabou's bicategory axioms use restricted amalgamated sums describable through multisimplices, suggesting a smaller starting category $\underline B_o$ than the generous one in the original letter. The handwritten cross-references on p.~16 revise the old page references to sections~8 and~9, and the scan preserves the tense and spelling ``This I ckecked'', where arXiv reads ``This I check''. The suggested nerve-like functor from $n$-stacks to $n$-semi-simplicial complexes is a goal for later investigation. Section~22 next reverses direction, seeking simpler special cases of categories, beginning with preorders and ordered sets. The uncertain final singular/plural in ``connection(s)'' on p.~15 remains an image-edge problem and has no mathematical consequence for this comparison.

\subsection*{22. Ordered sets, realizations and maquettes (pp. 16--19)}
Preorders and ordered sets offer simpler categorical presentations of homotopy types. Grothendieck relates ordered strata in conically stratified spaces to a realization whose barycentric simplices are finite totally ordered subsets, or flags. The resulting realization is canonically compared with that of the nerve. On p.~17 the corrected chain is $(\mathrm{Ord})\hookrightarrow(\mathrm{Preord})\hookrightarrow(\mathrm{Cat})\hookrightarrow(\mathrm{Ss\ sets})$, not the reversed first pair in arXiv. He cautions that the adjoint passage from a category to its object preorder can lose homotopy information, as the one-object case shows; the handwritten qualification ``but not Fl(A)'' is indispensable to that example. Page~18 gives a second realization through a nonseparated topology on the underlying ordered set and identifies sheaves with covariant functors, hence presheaves on the opposite order. Reversing the order preserves the flags and yields homeomorphic geometric realizations, while the two associated presheaf topoi need not be equivalent. The scan confirms hats in the latter claim and the author's rounded inclusion sign for the initial order comparison. Finally, maquettes are described by vertices and finite simplices and embedded in ordered sets through their face posets. The body claims that this category lacks finite products, while a dated marginal note adds that the notion is not an algebraic structure species; the hatching and call number of that note remain uncertain. Section~23 consequently shifts attention from satisfying Quillen's full package to constructing and localizing a natural functor into $(\mathrm{Hot})$. The edition preserves these claims and reservations as stages of the investigation, without treating the later witness's explanatory footnotes as text written on these pages.

\subsection*{23. The functor to homotopy types and a cellular core (pp. 19--21)}
Grothendieck separates the initial construction of a natural functor $M\to(\mathrm{Hot})$ from the further demand that $M$ carry a Quillen model structure. The functor would define weak equivalences by inverse image of isomorphisms; only then would one choose monomorphisms as cofibrations and test the resulting lifting-property definition of fibrations. The proposed site-theoretic construction uses the homotopy type of the induced site over each object, provided its pro-homotopy type is essentially constant. The canonical-topology suggestion is already qualified by a handwritten retraction pointing to section~24. Page~20 adds a structural requirement, universal and disjoint coproducts, and presents an algebraic species by a category $\underline C$ with finite inverse limits, its dual $\underline B$, and the anticipated equivalence $\mathrm{Ind}(\underline B)\simeq M$. This connects localization to finite presentations and to the question of which exactness properties are already visible in $\underline B$. Page~21 then isolates a smaller category $\underline B_{\mathrm{oo}}$ of basic cells and monomorphic boundary maps: the ordered set of proper subobjects should realize the boundary sphere of the corresponding cell. That construction uses section~22's ordered-set realizations and develops section~18's correction of the initial gluing intuition. The distinction between the broad species-level proposal and the geometric cell data remains heuristic. In the scan, the long p.~20 insertion includes ``can be interpreted'', omitted by arXiv, and the identity-functor arrow is drawn with two shafts; p.~21 preserves the printed slip ``of cubical'' and the typed replacements ``define'', ``family'' and ``sphere''.

\subsection*{24. Truncated nerve topologies and the failure of uniqueness (pp. 21--24)}
This section supplies the correction repeatedly anticipated in the earlier margins. Covering all composable pairs does not cover arbitrary strings: Grothendieck defines a descending sequence $T_N$ by epimorphic families of $N$-truncated nerves, identifying the first two levels with universal epimorphy and universal effective epimorphy. The comparison lemma relates the associated sheaf topoi to presheaves on truncated simplex categories. The finite levels retain only skeletal information, so the full nerve topology $T_\infty$ is needed for the proposed homotopy-type construction. The claim on p.~22 that an $N$-skeleton has the same homotopy and cohomology invariants through dimension $N$ requires caution at the endpoint: the $N$-skeleton of an $(N+1)$-simplex is an $N$-sphere, although the full simplex is contractible. The edition preserves the author's inequality rather than silently strengthening or repairing this intermediate assertion. The same loss of information explains the failure of $K\to NN'(K)$ discussed in section~16$'$: $N'$ factors through the two-dimensional data, while the nerve itself can still induce the desired equivalence after localization. Page~23 then replaces simplicial diagrams by commutative cubical diagrams, producing a strictly coarser topology with a different associated topos but the same intended homotopy interpretation. Multiple complexes suggest infinitely many further suitable topologies; the conclusion is that no unique topology is privileged. This leads directly to section~25's search for an equivalence relation among them. The scan's slice notation on p.~22 places $/A$ outside the parentheses ending after $T_N$, followed by an additional handwritten closing parenthesis; the earlier transcription's placement inside the subscript was corrected, while the extra parenthesis remains as written.

\subsection*{25. Hot-equivalence of topologies and global asphericity (pp. 24--25)}
The lack of a privileged suitable topology leads to a proposed equivalence relation. For comparable topologies $T\geq T'$ on $M$, Grothendieck compares the induced morphism of sheaf topoi and asks that every induced morphism over an object $A$ preserve homotopy type. He calls the resulting property ``globally aspheric'', strengthening the familiar cohomological notion of global acyclicity by allowing noncommutative coefficients. The primitive relation between comparable topologies is not yet an equivalence relation; taking the relation it generates is a first proposal, and he asks whether a coarser useful notion could do better. The intended invariant is a canonically associated functor $M\to(\mathrm{Hot})$, perhaps strong enough to recover the topology class itself. The analogy with Giraud's recovery of a topology from its associated subtopos motivates, rather than proves, a hoped-for correspondence. Page~25 narrows the aspiration to a single natural Hot-equivalence class for suitable species, admitting that even $M=(\mathrm{Cat})$ has no intrinsic characterization yet. Section~26 then reformulates the problem through small subcategories. Two revisions on p.~24 are essential: ``inverse image'' becomes ``direct image'', and the prime is moved so that the inclusion is from $T$-sheaves into $T'$-sheaves. The vertical margin note actually says ``just means'' and ``namely contractible'' with the qualification ``in the Čech sense''; the initial transcription's ``weakly contractible'' was not a literal reading. Its pencil shading leaves cancellation status uncertain. The body ends with the typed slip ``he'' for ``be'', and its final ``this'' is visibly canceled before the continuation on p.~25.

\subsection*{26. Test categories, test functors and the passage back to $M$ (pp. 25--29)}
The proposed intrinsic characterization of suitable topologies is reformulated through small categories $B$ inside a model category $M$. The cellular cores of section~23 and the simplex and cube categories of section~24 suggest that the presheaf category $\widehat B$ should itself describe homotopy types: its natural functor to $(\mathrm{Hot})$, computed from categories of objects over a presheaf, should be a localization. Grothendieck separates this internal requirement on $B$ from the eligibility of a functor $B\to M$; marginal terminology on p.~27 calls them ``test categories'' and ``test functors''. The comparison lemma gives an equivalence between the induced sheaf topos and $\widehat B$ under either of two fullness hypotheses. The decisive qualification comes on p.~28: localization of that sheaf category is not yet localization of $M$ itself. Generation by strict epimorphisms is considered as a way of making $M\to\widehat B$ fully faithful, but the author then proposes dropping both fullness requirements and directly demanding that $M\to\widetilde M\to\widehat B\to(\mathrm{Hot})$ be a localization. For categories $\underline B$ of finite presentations he further hopes for a localization onto homotopy types of finite type. The questions whether equal weak equivalences determine the functor, and whether suitable choices form one natural class, feed sections~27--28. The margin of p.~27 actually says ``same question arises'' for Hot-equivalence of topologies; ``arises'' was absent from the initial transcription. Its wording is now readable, but pencil hatching still leaves the note's cancellation status uncertain. The call after objects ``over a'' on p.~26 is (23), not (3).

\subsection*{27. Geometric realization as a third construction (p. 29)}
Calling the preceding discussion ``thin air conjecturing'', Grothendieck warns that both test categories and test functors may need stronger conditions. He places geometric or ``topologic'' realization beside the two constructions through topologies and through test functors. Semi-simplicial and cubical complexes offer comparatively direct realizations, hence so do the categories embedded in semi-simplicial sets. For Gr-stacks, the prescribed realization of the primitive $\underline B_{\mathrm{oo}}$ structure has to extend to a left coherator, and that extension involves the arbitrary choices already encountered in sections~18--19. Rather than immediately eliminating them, he compares this indeterminacy with the variety of test topologies and functors developed in sections~24--26. The paragraph thus states the next intended problem, relating realization to the two test constructions; section~28 immediately postpones it in favor of rigidity of $(\mathrm{Hot})$. The scan preserves ``topologic'', where the witness expands it to ``topological'', and the word ``is'' in ``is that a suitable'', where the witness has ``it''. Its call (28) supplies no text of the witness's retrospective footnote. The typed insertions ``localization'', ``left'' and ``between'' sharpen the objects of the comparison; the two round letters in $\underline B_{\mathrm{oo}}$ and the primitive graph designation remain literal rather than silently replaced by a joined infinity glyph.

\subsection*{28. Rigidity of $(\mathrm{Hot})$ and the first modelizers (pp. 30--33)}
The intended return to geometric realization is suspended by an ``Assumption'': the category of self-equivalences of $(\mathrm{Hot})$ and their natural isomorphisms should be equivalent to a point. Grothendieck separates its two demands, every self-equivalence is isomorphic to the identity and the identity has no nontrivial automorphism, while wondering whether arbitrary natural endomorphisms could also be admitted. This is asserted as a plausible assumption, not proved; it supplies the desired uniqueness up to unique isomorphism of a localization functor once its weak equivalences are known. A ``modelizing category'' or ``modelizer'' is then a category $M$ with a saturated class $W$ whose localization is equivalent to $(\mathrm{Hot})$. A model-preserving functor carries weak equivalences to weak equivalences and induces an equivalence of localizations, equivalently commuting up to isomorphism with the type functors. This formalizes sections~26--27's questions without yet specifying cofibrations and fibrations. The diagram on p.~31 compares ordered sets, preorders, categories, semi-simplicial sets and cubical sets; the category-of-elements functors $i_A$ lead to right adjoints $j_A$. Page~32 carefully distinguishes those right adjoints from the familiar nerve functors, contrasting infinite categories of simplices with finite ordered face structures when degeneracies are omitted. Mixed multicomplexes then enlarge the diagram and motivate the test-category definition in section~29. The mixed-product display on p.~32 visibly has $\Delta$ in both factors, although the adjacent prose says $p$ times simplicial and $q$ times cubical; this apparent authorial symbol slip is preserved. The same page changes ``modelators'' to ``modelizers'' by hand, a correction the witness does not apply, and contains a small branching diagram omitted there. The p.~30 cross-reference is to section~26; calls (28$'$), (29) and (30) are present without the witness's expanded footnotes.

\subsection*{29. The basic modelizer and a stronger test condition (pp. 33--36)}
The category of small categories now becomes the foundational modelizer: $(\mathrm{Hot})$ is its localization, and weak equivalences are described through the induced morphisms of presheaf topoi and the Artin--Mazur cohomological criterion recalled from section~17. A small category $A$ is provisionally a test category when $i_A\colon\widehat A\to(\mathrm{Cat})$, sending a presheaf to its category of elements, is model preserving; its presheaf category is an ``elementary modelizer''. This replaces a special preference for semi-simplicial sets by a comparison among all such presentations, while acknowledging the computational convenience of the nerve and the simplicial construction of pro-homotopy types. The general adjunction calculation of pp.~34--35 asks when both adjoints induce inverse equivalences after localization. Applied to $i_A\dashv j_A$, it leads to condition (T): every counit $i_Aj_A(B)\to B$ should be a weak equivalence. Grothendieck explicitly leaves open whether this strengthens the preceding definition, and proposes strengthening it if necessary; the call (33) does not contain the witness's retrospective answer. Products of test categories and direct comparison of elementary models become the next questions. The closing heuristic on p.~36 would test (T) by a fibration with contractible fibers, either after taking nerves or directly in $(\mathrm{Cat})$; section~30 immediately attacks the model-category assumption behind the latter route. On p.~35 the source twice writes that $M'$ or $M$ is saturated, although the discussion concerns the arrow classes $W'$ and $W$; these apparent symbol slips remain unchanged. The typed ``elementary normalizer'' on that page likewise remains source text. The scan's baseline slash in $M/_{b}$ on p.~34 is preserved, and its reference reads ``section 28'', not the witness's page reference. The broad emphasis of the witness also differs from the selective underlining and marginal double bar of p.~33.

\subsection*{30. From the model-category obstruction to asphericity (pp. 36--38)}
The attempt to treat every monomorphism of categories as a cofibration now runs into the two elementary horns drawn on p.~37. Lifting against both would force the arrows of a fibrant category to have left and right inverses, leaving groupoids, whose homotopy types cannot exhaust those of arbitrary categories. The objection concerns that proposed choice of cofibrations; the source does not prove the nonexistence of every possible model structure on categories. Grothendieck nevertheless keeps the expectation that elementary modelizers should admit closed model structures, and turns instead to the asphericity criterion introduced in section~25. A morphism of topoi is required to remain a weak equivalence after every localization over the target; when the topoi come from categories, generators reduce the condition to the asphericity of each slice $C'_{/a}$. Applying this to the counit $i_Aj_A(B)\to B$ from section~29 gives three equivalent conditions: the strengthened test condition, asphericity of every counit, and asphericity of $i_Aj_A(B)$ whenever $B$ has a final element. That last formulation opens the explicit interval argument of section~31. The scan's four bracketed ``[small]'' restrictions on p.~38 are genuine hand additions, including two absent from the witness. Its caret under the $A$ in $\operatorname{Hom}(A_{/a},B)$ is retained in its actual position without asserting a mathematical interpretation. The apparent inline hat on $f$ reported by the initial comparison on p.~37 belongs instead to the display beneath it; the inline symbol has been corrected to plain $f$. The word ``closed'' on p.~37 and the replacement ``aspheric'' on p.~38 are interlinear typing, although their placing strokes are handwritten.
